\begin{explanationbox}
	\textbf{\textcolor{CExplanation}{一句话直觉}}

	把自变量 $x$ 换成 $-x$：奇函数的函数值变号，偶函数的函数值不变。
	看运算后的整个式子是否变号，就能判断结果的奇偶性。

	\medskip
	\textbf{\textcolor{CExplanation}{加减法：看能否整体提出符号}}

	如果 $f$、$g$ 都是奇函数，那么
	\[
		f(-x)+g(-x)=-f(x)-g(x)=-\bigl(f(x)+g(x)\bigr).
	\]
	整个和变号，所以仍为奇函数。同为偶函数时，两项都不变，
	和或差也不变。

	若一个奇、一个偶，例如 $x+x^2$，代入 $-x$ 后只有其中一项变号，
	通常无法给整个式子提出同一个正号或负号。

	\medskip
	\textbf{\textcolor{CExplanation}{乘除法：数变号的次数}}

	奇函数因子遇到 $-x$ 会变号，偶函数因子不会。
	两个奇函数相乘，变号两次，结果不变，所以乘积为偶函数；
	只有一个奇函数因子时，结果变号，所以乘积为奇函数。

	商也按同样方法判断，但先要删去分母等于零的点。
	例如 $x/x$ 在 $x\ne 0$ 上等于 $1$，是偶函数；
	不能因为它约去后等于 $1$，就把 $x=0$ 加回原定义域。

	\medskip
	\textbf{\textcolor{CExplanation}{记忆方法}}

	乘除法可记为：\textbf{出现奇函数的次数为奇数次，结果为奇；
		出现偶数次，结果为偶。} 加减法则要分别检查每一项，
	不能直接套用这条口诀。
\end{explanationbox}